\chapter{Results \& Analysis}
\label{ch:results}

\section{Baseline Performance: Inductive Conformal Prediction}

The ICP framework was configured with a target error rate of \(\alpha = 0.1\),
corresponding to a nominal coverage level of 90\% for each of the 14 pathologies.

Figure~\ref{fig:baseline-coverage} presents the achieved coverage rate
for each pathology on the NIH \gls{cxr} test set,
using the DenseNet121 model.
The y-axis represents the empirical coverage rate,
calculated as the proportion of test samples for a given pathology
where the true label was correctly included in the prediction set.
The x-axis lists each of the 14 distinct pathologies evaluated.
The dashed red line indicates the nominal coverage target of 90\% \((1 - \alpha)\).

\begin{figure}
    \centering
    \includegraphics{figures/results-baseline-coverage}
    \caption{%
        Achieved class-conditional coverage rates, using DenseNet121 on the NIH dataset.
        The results confirm the validity of the implementation,
        as the empirical coverage for every class meets or exceeds
        the 90\% nominal level (red line).
    }
    \label{fig:baseline-coverage}
\end{figure}

The results depicted clearly validate the correctness
of the class-conditional ICP implementation.
For every pathology, the empirical coverage rate meets or exceeds the 90\% target.
The variation in coverage, particularly the high rate for Pneumothorax,
suggests that the model has differing levels of difficulty for each class,
forcing the prediction set to be more conservative for some than for others.

While ICP provides valid coverage,
its primary objective is not to control specific types of errors.
The critical question for clinical deployment is whether it is suitable
for risk-sensitive applications, where the cost of false negatives
is significantly higher than that of false positives.
To investigate this, the FNR of ICP was compared against that of CRC,
which was configured to target an FNR of \(\alpha=0.1\).
Both frameworks were evaluated under 100 different random test splits,
with the distribution of FNRs shown in Figure~\ref{fig:baseline-crc-risk}.

\begin{figure}
    \centering
    \begin{subfigure}[b]{0.40\textwidth}
        \centering
        \includegraphics[width=\linewidth]{figures/results-baseline-crc-risk}
        \caption{%
            Distribution of FNRs for ICP and CRC.
        }
        \label{fig:baseline-crc-risk}
    \end{subfigure}
    \hfill
    \begin{subfigure}{0.56\textwidth}
        \centering
        \includegraphics[width=\linewidth]{figures/results-baseline-crc-size}
        \caption{%
            Distribution of prediction set sizes for ICP and CRC.
        }
        \label{fig:baseline-crc-size}
    \end{subfigure}
    \caption{%
        \textcolor{red}{(TODO: legend overlaps, y-axis needs to be consistent with other fig)}
        Comparison of ICP and CRC in terms of FNR and prediction set size,
        using DenseNet121 on the NIH dataset.
        (a) shows the distribution of FNRs, with the dashed line
        indicating the target level of 10\% for CRC.
        (b) illustrates the difference in distribution of prediction set sizes.
    }
\end{figure}

The result is unambiguous. The FNR for ICP is high,
centred at approximately 0.43.
In stark contrast, the FNR for CRC is tightly controlled,
with the majority of its distribution falling below the 10\% target level.
This empirically proves that CRC provides effective and reliable control of the FNR,
while standard ICP does not.
It is noted that ICP was not expected to control FNR,
but this result highlights the limitations of ICP
when applied to risk-sensitive applications.

To understand the practical implications of these different objectives,
Figure~\ref{fig:baseline-crc-size} compares the distribution of
prediction set sizes produced by each method.
A significant number of predictions are the trivial set of all 14 labels,
indicating catastrophic model uncertainty.
The CRC framework, while also favouring small sets,
almost entirely avoids producing these large, clinically uninformative sets.
This finding suggests that by focusing on a specific risk (FNR),
CRC not only achieves a safer outcome
but also does so with greater overall predictive efficiency,
as it is not forced into the same `all-or-nothing' predictions as the baseline.


\textcolor{red}{TODO: Get example image from dataset and both prediction sets.}

\section{Monotone Risk Control}

By defining and implementing a custom,
clinically-weighted loss function that assigns a higher penalty to missed pathologies,
the aim is to evaluate whether CRC can be tailored to better align with clinical priorities.
It is important to note that the author is not a clinician,
and these weights do not constitute a clinical standard.
Rather, they serve as a proof-of-concept to demonstrate the capabilities of the methodology.

\textcolor{red}{TODO: Graph and conclusion.}

\section{High-Probability Guarantees with LTT}

This section provides a direct, empirical comparison between CRC and LTT,
first by evaluating the reliability of their guarantees on the same monotone risk (FNR) and second,
by demonstrating LTT's unique ability to control non-monotone risks that are beyond the scope of CRC.

To conduct a direct comparison of the guarantees,
both CRC and LTT were tasked with controlling the FNR at a target level of \(\alpha = 0.1\).
For the LTT framework, the high-probability parameter was set to \(\delta = 0.1\),
meaning it is expected to control the risk with at least 90\% probability.

\begin{figure}
    \centering
    \begin{subfigure}[b]{0.40\textwidth}
        \centering
        \includegraphics[width=\linewidth]{figures/results-crc-ltt-risk}
        \caption{%
            Distribution of FNRs for CRC and LTT,
            with both frameworks targeting an FNR of 10\%.
        }
        \label{fig:results-crc-ltt-risk}
    \end{subfigure}
    \hfill
    \begin{subfigure}{0.56\textwidth}
        \centering
        \includegraphics[width=\linewidth]{figures/results-crc-ltt-size}
        \caption{%
            Distribution of prediction set sizes for CRC and LTT.
        }
        \label{fig:crc-ltt-size}
    \end{subfigure}
    \caption{%
        Comparison of CRC and LTT in terms of FNR and prediction set size,
        using DenseNet121 on the NIH dataset.
    }
\end{figure}

Figure~\ref{fig:results-crc-ltt-risk} presents the distribution of
the achieved FNR for both frameworks over 100 random test splits.
The distribution for CRC is centred, as expected, near the nominal 0.1 level.
However, it exhibits significant variance, with a substantial portion of the trials
resulting in an FNR that exceeds the target \(\alpha\).
This empirically demonstrates the nature of an in-expectation guarantee:
while the average performance is controlled,
there is no protection against underperformance on any given trial.
For safety-critical systems, these violations can represent unacceptable failures.

In stark contrast, the LTT framework demonstrates a fundamentally more robust behaviour.
The distribution of its achieved FNR, as shown in the same figure,
is sharply truncated at the \(\alpha\) threshold.
In over 90\% of the trials, the risk was held at or below the target level.
This is the practical manifestation of the \((1 - \delta)\) high-probability guarantee:
LTT provides a reliable `risk ceiling',
ensuring that the system's performance will not degrade beyond the specified level,
except in a small, user-defined fraction \(\delta\) of cases.

Naturally, this enhanced reliability is not without cost.
Figure~\ref{fig:crc-ltt-size} presents the corresponding distribution of prediction set sizes.
To achieve its stronger guarantee,
LTT produces prediction sets that are slightly larger on average than those from CRC.
This represents a trade-off:
a small decrease in the precision of the prediction in exchange for a significant increase
in its reliability.

\textcolor{red}{TODO: Add example image from dataset and prediction sets.}

% To ground these statistical distributions in a clinical context,
% a case study was selected from a calibration split where CRC's performance was notably poor,
% resulting in an achieved FNR of 0.18.
% Figure 5.9 illustrates a scenario corresponding to this challenging calibration split.
% The radiograph presents a subtle case of Pleural Effusion
% that is difficult for the model to detect with high confidence.
% For this `unlucky' split, the calibration data was not perfectly representative
% of the model's general performance,
% causing the CRC algorithm to adopt an overly optimistic threshold \(\lambda\).
% Consequently, when presented with the test case, its prediction set was empty,
% thereby missing the critical pathology.

% The LTT framework, when run on the exact same data split,
% was insulated from this instability.
% It selected a safer, more conservative threshold
% and correctly included 'Pleural Effusion' in its prediction set.
% This case highlights how LTT's high-probability guarantee
% provides a crucial layer of safety against unfavourable data splits,
% preventing model failure in scenarios where a less robust framework would have failed.

Beyond offering a stronger guarantee on monotone risks,
the primary advantage of the LTT framework is its ability to control arbitrary risk functions
that are non-monotone;
a capability fundamentally beyond the scope of CRC.

\begin{figure}
    \centering
    \includegraphics{figures/results-ltt-risk-control}
    \caption{%
        Achieved FDR on the test set, using LTT to control the risk at various target levels.
        The results demonstrate LTT's ability to manage non-monotone risks effectively.
    }
    \label{fig:results-ltt-fdr}
\end{figure}

Figure~\ref{fig:results-ltt-fdr} presents the results of applying LTT
to control FDR at various target levels, from \(\alpha = 0.05 \text{ to } \alpha = 0.5\),
over 100 random test splits.
As clearly shown, the achieved FDR on the test set
tracks consistently below the nominal target \(\alpha\).
This provides strong evidence that the LTT framework
successfully controls the FDR at the desired level,
and the fact that the achieved risk does not substantially undercut the target
indicates that the procedure is not overly conservative.
It is efficient, providing a level of risk control that is tight
to the user-specified boundary without being unnecessarily restrictive,
which would otherwise lead to overly small (and less useful) prediction sets.

This result is of practical importance, confirming that by
reframing calibration as a multiple hypothesis testing problem,
LTT successfully circumvents the monotonicity limitation inherent to CRC.
This empowers practitioners to manage a different, yet equally critical, clinical risk:
the danger of `alarm fatigue',
where a model that makes too many false positive predictions
may lead clinicians to distrust its outputs.
By allowing for the control of risks like FNR and FDR,
LTT provides a far more comprehensive proceedure
for building truly reliable and useful clinical AI systems.

\section{Robustness and Generalisation}

The previous sections have established the comparative performance of ICP, CRC,
and LTT on the NIH dataset using a DenseNet121 architecture.
However, for these findings to be truly meaningful,
they must be robust to changes in the underlying model and the data distribution.
This section, therefore, aims to validate the generalizability of results
by systematically evaluating the frameworks across two distinct datasets
(CheXpert and NIH ChestX-ray14)
and two different model architectures (DenseNet121 and ResNet50).
